\documentclass[11pt,a4paper]{report}

% ─────────────────────────────────────────────
%  Packages
% ─────────────────────────────────────────────
\usepackage[margin=2.5cm]{geometry}
\usepackage[T1]{fontenc}
\usepackage[utf8]{inputenc}
\usepackage{babel}

\usepackage{graphicx}
\usepackage{hyperref}
\usepackage{titlesec}
\usepackage{fancyhdr}
\usepackage{booktabs}
\usepackage{tabularx}
\usepackage{longtable}
\usepackage{array}
\usepackage{listings}
\usepackage{xcolor}
\usepackage{setspace}
\usepackage{parskip}

\usepackage{enumitem}
\usepackage{amsmath}
\usepackage{caption}
\usepackage{tocloft}
\usepackage{emptypage}

% ─────────────────────────────────────────────
%  Hyperref setup
% ─────────────────────────────────────────────
\hypersetup{
    colorlinks=true,
    linkcolor=blue!60!black,
    urlcolor=blue!60!black,
    citecolor=blue!60!black,
    pdftitle={Système de Pilotage de la Performance Organisationnelle des Interventions},
    pdfauthor={Yassine CHELLY \& Amen Allah BEZZAOUIA},
}

% ─────────────────────────────────────────────
%  Code listings style
% ─────────────────────────────────────────────
\lstdefinestyle{jsStyle}{
    language=Java,
    basicstyle=\ttfamily\small,
    keywordstyle=\color{blue!70!black}\bfseries,
    stringstyle=\color{orange!80!black},
    commentstyle=\color{gray}\itshape,
    breaklines=true,
    frame=single,
    backgroundcolor=\color{gray!10},
    rulecolor=\color{gray!40},
    numbers=left,
    numberstyle=\tiny\color{gray},
    tabsize=2,
    showstringspaces=false,
    captionpos=b,
}

% ─────────────────────────────────────────────
%  Section formatting
% ─────────────────────────────────────────────
\titleformat{\chapter}[display]
  {\normalfont\huge\bfseries\color{blue!60!black}}
  {\chaptertitlename\ \thechapter}{16pt}{\Huge}

\titleformat{\section}
  {\normalfont\Large\bfseries\color{blue!50!black}}
  {\thesection}{1em}{}

\titleformat{\subsection}
  {\normalfont\large\bfseries}
  {\thesubsection}{1em}{}

\titleformat{\subsubsection}
  {\normalfont\normalsize\bfseries}
  {\thesubsubsection}{1em}{}

% ─────────────────────────────────────────────
%  Header / footer
% ─────────────────────────────────────────────
\pagestyle{fancy}
\fancyhf{}
\rhead{\small SPPOI -- Tunisie Telecom}
\lhead{\small ENICarthage 2025--2026}
\cfoot{\thepage}
\renewcommand{\headrulewidth}{0.4pt}

% ─────────────────────────────────────────────
%  Line spacing
% ─────────────────────────────────────────────
\setstretch{1.2}

% ─────────────────────────────────────────────
%  Column type for tables
% ─────────────────────────────────────────────
\newcolumntype{L}[1]{>{\raggedright\arraybackslash}p{#1}}
\newcolumntype{C}[1]{>{\centering\arraybackslash}p{#1}}

% ═════════════════════════════════════════════
%  DOCUMENT
% ═════════════════════════════════════════════
\begin{document}

% ─────────────────────────────────────────────
%  PAGE DE TITRE
% ─────────────────────────────────────────────
\begin{titlepage}
\centering
\vspace*{1cm}

{\large\bfseries ÉCOLE NATIONALE D'INGÉNIEURS DE CARTHAGE (ENICarthage)}\\[0.3cm]
{\large Filière : Génie des Systèmes Industriels et Logistiques (ISL)}\\[0.3cm]
{\large Année universitaire 2025--2026}

\vspace{1.5cm}
\rule{\linewidth}{1.5pt}\\[0.4cm]

{\Huge\bfseries\color{blue!60!black}
Système de Pilotage de la Performance Organisationnelle des Interventions}

\vspace{0.4cm}
\rule{\linewidth}{1.5pt}

\vspace{1cm}
{\LARGE\bfseries RAPPORT DE PROJET DE FIN D'ANNÉE}

\vspace{1cm}
{\large Entreprise d'accueil : \textbf{Tunisie Telecom}}

\vfill

\begin{minipage}[t]{0.45\textwidth}
\raggedright
\textbf{Réalisé par :}\\[0.3cm]
Yassine \textsc{Chelly}\\
Amen Allah \textsc{Bezzaouia}
\end{minipage}
\hfill
\begin{minipage}[t]{0.45\textwidth}
\raggedleft
\textbf{Encadrante :}\\[0.3cm]
Mme Haifa \textsc{Grira}
\end{minipage}

\vspace{1.5cm}
{\large Année 2025--2026}
\end{titlepage}

% ─────────────────────────────────────────────
%  RÉSUMÉ
% ─────────────────────────────────────────────
\chapter*{Résumé}
\addcontentsline{toc}{chapter}{Résumé}

Dans un contexte de transformation digitale accélérée du secteur des télécommunications en Tunisie, la gestion efficace des réclamations et des interventions techniques constitue un enjeu stratégique majeur pour les opérateurs. Ce projet de fin d'année, réalisé au sein de Tunisie Telecom en collaboration avec le fournisseur de services Internet Hexabyte, vise à concevoir et implémenter un Système de Pilotage de la Performance Organisationnelle des Interventions (SPPOI).

Face au constat d'une absence d'outil centralisé de pilotage et d'une visibilité parcellaire sur les 22~851 réclamations traitées sur l'année 2025 (taux de résolution~: 99,93~\%, délai moyen~: 79 heures, 23,1~\% de tickets urgents dépassant 72 heures), le système développé repose sur une architecture applicative multi-interfaces en Next.js, déployée sur Vercel. Il intègre un dashboard de performance principale offrant sept vues analytiques interactives, un module d'analyse Pareto intelligente exploitant l'API Groq et le modèle de langage Llama~3.3~70B pour la génération automatisée de rapports de diagnostic, ainsi que des interfaces dédiées aux fournisseurs FSI et aux agents techniciens de Tunisie Telecom.

\medskip
\noindent\textbf{Mots-clés~:} Pilotage de la performance, KPI, Tableau de bord, Pareto intelligent, Next.js, Intelligence artificielle, Tunisie Telecom, Gestion des interventions, LLM, Groq.

% ─────────────────────────────────────────────
%  ABSTRACT
% ─────────────────────────────────────────────
\chapter*{Abstract}
\addcontentsline{toc}{chapter}{Abstract}

In a context of accelerated digital transformation within Tunisia's telecommunications sector, the effective management of technical complaints and field interventions represents a major strategic challenge for operators. This final year project, carried out within Tunisie Telecom in collaboration with Internet service provider Hexabyte, aims to design and implement an Organizational Performance Monitoring System for Technical Interventions (OPMSTI).

Faced with the absence of a centralized monitoring tool and fragmented visibility over the 22,851 complaints processed in 2025 (resolution rate: 99.93\%, average processing time: 79 hours, 23.1\% of urgent tickets exceeding 72 hours), the developed system is built on a multi-interface Next.js application architecture deployed on Vercel.

\medskip
\noindent\textbf{Keywords:} Performance monitoring, KPI, Dashboard, Intelligent Pareto, Next.js, Artificial intelligence, Tunisie Telecom, Intervention management, LLM, Groq.

% ─────────────────────────────────────────────
%  REMERCIEMENTS
% ─────────────────────────────────────────────
\chapter*{Remerciements}
\addcontentsline{toc}{chapter}{Remerciements}

Nous tenons, en premier lieu, à exprimer notre profonde gratitude envers notre encadrante académique, Madame Haifa \textsc{Grira}, dont les conseils avisés, la disponibilité constante et la rigueur intellectuelle ont guidé chacune des étapes de ce projet. Sa bienveillance et son exigence ont été pour nous une source de motivation tout au long de cette année universitaire.

Nous adressons également nos sincères remerciements à l'ensemble des équipes de Tunisie Telecom qui nous ont accueillis au sein de leur direction des opérations, et plus particulièrement aux cadres et techniciens qui ont accepté de partager leur expertise opérationnelle et de nous fournir les données nécessaires à la réalisation de ce projet.

Nous exprimons notre reconnaissance à l'École Nationale d'Ingénieurs de Carthage (ENICarthage) et à l'ensemble du corps enseignant de la filière Génie des Systèmes Industriels et Logistiques, pour la qualité de la formation dispensée et pour nous avoir préparés à relever les défis d'un projet de cette envergure.

Nos remerciements vont également aux membres du jury, qui nous font l'honneur d'évaluer ce travail. Leurs retours et observations contribueront sans nul doute à enrichir notre réflexion et à renforcer la qualité de nos futurs travaux.

Enfin, nous souhaitons exprimer notre reconnaissance à nos familles respectives, pour leur soutien indéfectible, leur patience et leurs encouragements tout au long de ce parcours académique et professionnel.

% ─────────────────────────────────────────────
%  LISTE DES ABRÉVIATIONS
% ─────────────────────────────────────────────
\chapter*{Liste des Abréviations}
\addcontentsline{toc}{chapter}{Liste des Abréviations}

\begin{longtable}{L{3cm} L{11cm}}
\toprule
\textbf{Abréviation} & \textbf{Signification} \\
\midrule
\endhead
\bottomrule
\endfoot
ADSL    & Asymmetric Digital Subscriber Line -- Ligne numérique asymétrique à paires torsadées \\
API     & Application Programming Interface -- Interface de programmation applicative \\
CI/CD   & Continuous Integration / Continuous Deployment \\
CSS     & Cascading Style Sheets -- Feuilles de style en cascade \\
ENICarthage & École Nationale d'Ingénieurs de Carthage \\
FSI     & Fournisseur de Services Internet \\
HPSM    & HP Service Manager -- Système de gestion des incidents de Tunisie Telecom \\
IA      & Intelligence Artificielle \\
ISL     & Génie des Systèmes Industriels et Logistiques \\
KPI     & Key Performance Indicator -- Indicateur Clé de Performance \\
LLM     & Large Language Model -- Grand modèle de langage \\
OPMSTI  & Organizational Performance Monitoring System for Technical Interventions \\
PFA     & Projet de Fin d'Année \\
React   & Bibliothèque JavaScript open-source pour la création d'interfaces utilisateur \\
SLA     & Service Level Agreement -- Accord de niveau de service \\
SMART   & Specific, Measurable, Achievable, Realistic, Time-bound \\
SPPOI   & Système de Pilotage de la Performance Organisationnelle des Interventions \\
SSG     & Static Site Generation -- Génération de site statique (Next.js) \\
SSR     & Server-Side Rendering -- Rendu côté serveur (Next.js) \\
TT      & Tunisie Telecom \\
UML     & Unified Modeling Language -- Langage de modélisation unifié \\
VDSL    & Very High Speed Digital Subscriber Line -- Ligne numérique à très haut débit \\
\end{longtable}

% ─────────────────────────────────────────────
%  TABLE DES MATIÈRES
% ─────────────────────────────────────────────
\tableofcontents
\listoftables

\clearpage

% ═════════════════════════════════════════════
%  INTRODUCTION GÉNÉRALE
% ═════════════════════════════════════════════
\chapter*{Introduction Générale}
\addcontentsline{toc}{chapter}{Introduction Générale}

\section*{Contexte général}

Le secteur des télécommunications constitue l'une des industries les plus stratégiques et les plus dynamiques de l'économie tunisienne. Dans un contexte mondial marqué par une accélération sans précédent de la transformation digitale, les opérateurs télécom font face à des exigences croissantes en matière de qualité de service, de réactivité opérationnelle et de pilotage de la performance. Tunisie Telecom, opérateur historique et principal acteur du marché national, assure la gestion d'un réseau d'accès fixe couvrant l'ensemble du territoire national, dont la qualité repose en grande partie sur la fluidité du traitement des réclamations techniques de ses clients.

La gestion des réclamations et des interventions techniques constitue le c\oe ur de l'activité opérationnelle de Tunisie Telecom. Chaque incident client -- qu'il s'agisse d'une perte de synchronisation, d'une déconnexion courante ou d'une dégradation du débit -- génère un ticket d'intervention qui doit être pris en charge dans les meilleurs délais par un agent technique. La multiplicité des acteurs impliqués -- fournisseurs de services Internet (FSI), agents terrain de Tunisie Telecom, superviseurs régionaux -- ainsi que la diversité géographique et technique des interventions rendent le pilotage de cette activité particulièrement complexe.

\section*{Problématique}

L'analyse du système existant au sein de Tunisie Telecom révèle plusieurs lacunes opérationnelles significatives. En l'absence d'un outil centralisé de pilotage, les responsables des opérations se trouvent confrontés à une visibilité parcellaire sur l'état réel des interventions en cours et sur les performances des différents acteurs de la chaîne d'intervention.

La question centrale à laquelle ce projet entend répondre est donc la suivante :

\begin{quote}
\itshape Comment concevoir et implémenter un système de pilotage de la performance organisationnelle des interventions permettant aux responsables de Tunisie Telecom de disposer, à tout moment, d'une vision claire et actionnable des indicateurs clés de performance, d'identifier les principales causes de dysfonctionnement, et d'initier des actions correctives ciblées~?
\end{quote}

\section*{Objectifs du projet}

En réponse à la problématique identifiée, ce projet de fin d'année s'est fixé les quatre objectifs principaux suivants :

\begin{enumerate}[leftmargin=*]
  \item Développer un système de pilotage de la performance organisationnelle intégrant les données réelles issues des fournisseurs et des agents Tunisie Telecom, et restituant des indicateurs clés de performance sous forme de tableaux de bord interactifs et dynamiques.
  \item Mettre en \oe uvre une analyse Pareto intelligente, assistée par l'intelligence artificielle, permettant d'identifier automatiquement les 20~\% de causes générant 80~\% des incidents, et de générer des rapports de diagnostic contextualisés.
  \item Concevoir une architecture applicative multi-interfaces adaptée aux besoins spécifiques de chaque acteur~: interface FSI, interface agent, et dashboard de performance.
  \item Fournir un outil d'aide à la décision robuste, basé sur des données réelles, permettant aux décideurs de piloter efficacement la performance organisationnelle.
\end{enumerate}

\section*{Organisation du rapport}

Le présent rapport est organisé en cinq chapitres complémentaires~:
\begin{itemize}
  \item \textbf{Chapitre 1}~: Présentation de l'organisme d'accueil et cadre du projet.
  \item \textbf{Chapitre 2}~: État de l'art et cadre théorique.
  \item \textbf{Chapitre 3}~: Analyse et conception du système.
  \item \textbf{Chapitre 4}~: Réalisation et implémentation.
  \item \textbf{Chapitre 5}~: Tests, validation et résultats.
\end{itemize}

% ═════════════════════════════════════════════
%  CHAPITRE 1
% ═════════════════════════════════════════════
\chapter{Présentation de l'Organisme d'Accueil et Cadre du Projet}

Ce premier chapitre a pour objet de situer le projet dans son environnement institutionnel et opérationnel. Nous présentons successivement l'organisme d'accueil Tunisie Telecom, le contexte et le diagnostic du système existant, puis nous formulons le cahier des charges du système à développer.

\section{Présentation de Tunisie Telecom}

\subsection{Historique et évolution}

Tunisie Telecom est l'opérateur public de télécommunications de la Tunisie. Fondée en 1995 sous sa forme actuelle, suite à la restructuration de la direction générale des télécommunications relevant du Ministère des Technologies de la Communication, Tunisie Telecom a connu une évolution continue pour s'affirmer comme le principal opérateur multiservices du pays. Depuis ses débuts en tant que gestionnaire exclusif du réseau téléphonique commuté (RTC), l'opérateur a progressivement étendu son portefeuille de services vers le haut débit fixe (ADSL, VDSL), le mobile (2G, 3G, 4G et 4G+), et les services aux entreprises.

L'ouverture progressive du marché tunisien des télécommunications à la concurrence, avec l'entrée d'Ooredoo en 2002 et d'Orange Tunisie en 2010, a contraint Tunisie Telecom à accélérer sa modernisation et à investir massivement dans la qualité de service afin de maintenir sa position de leader.

\subsection{Missions, vision et orientations stratégiques}

La mission principale de Tunisie Telecom est de fournir à l'ensemble des citoyens tunisiens un accès universel et de qualité aux services de télécommunications. L'opérateur s'est doté d'une vision stratégique centrée sur trois axes majeurs~: la modernisation des infrastructures réseaux, l'amélioration continue de l'expérience client, et le développement de nouveaux services à valeur ajoutée.

\subsection{Structure organisationnelle}

Tunisie Telecom est organisée en directions fonctionnelles et opérationnelles. La direction des opérations, au sein de laquelle s'inscrit ce projet, est responsable de la gestion du réseau d'accès fixe, de la supervision des interventions techniques et de la qualité de service délivrée aux clients finaux.

\subsection{Portefeuille de services}

Tunisie Telecom propose une gamme étendue de services couvrant les besoins des particuliers et des entreprises. Sur le segment fixe, l'opérateur commercialise des offres ADSL et VDSL sous diverses formules. L'Offre de Gros ADSL Hybride représente à elle seule 66,3~\% du volume total des réclamations traitées.

\subsection{Position concurrentielle sur le marché télécom tunisien}

Malgré la concurrence accrue, Tunisie Telecom maintient une position dominante sur le marché de l'accès fixe à internet haut débit, grâce à l'étendue de son réseau de cuivre et à son déploiement progressif de la fibre optique.

\section{Contexte et Diagnostic du Système Existant}

\subsection{Processus actuel de gestion des réclamations et interventions}

Le processus de gestion des réclamations et interventions chez Tunisie Telecom implique plusieurs acteurs et systèmes d'information distincts. Lorsqu'un client signale une anomalie sur sa ligne, le FSI -- en l'occurrence Hexabyte -- effectue un premier diagnostic à distance. Si l'anomalie ne peut être résolue à distance, le FSI ouvre un ticket de réclamation sur le portail Extranet Tunisie Telecom.

Le ticket est ensuite pris en charge par le système HPSM (HP Service Manager) de Tunisie Telecom, qui lui attribue une référence unique et le route vers l'équipe d'intervention compétente. Un agent terrain réalise l'intervention physique et renseigne le code de résolution ainsi que la famille d'intervention.

\subsection{Identification des dysfonctionnements et limites du système existant}

L'analyse des données réelles disponibles, portant sur 22~851 tickets traités sur l'année 2025, met en évidence plusieurs dysfonctionnements structurels~:

\begin{itemize}
  \item \textbf{Absence de visibilité en temps réel}~: les responsables ne disposent d'aucun tableau de bord consolidé permettant de suivre en temps réel le volume des tickets ouverts, les délais de traitement et les taux de résolution.
  \item \textbf{Délais de traitement élevés et hétérogènes}~: si le taux de résolution global est excellent (99,9~\%), le délai moyen de traitement atteint 79 heures, avec des disparités régionales importantes (Mannouba~: 180,8~h, contre 21,2~h pour les interventions classées "Rien à signaler").
  \item \textbf{Concentration des incidents non identifiée}~: l'analyse Pareto révèle que six codes de résolution concentrent plus de 81,6~\% des réclamations.
  \item \textbf{Tickets urgents en volume significatif}~: 23,1~\% des tickets (5~278 tickets) dépassent le seuil critique de 72 heures sans résolution.
\end{itemize}

\section{Cahier des Charges}

\subsection{Besoins fonctionnels}

Sur la base du diagnostic établi, les besoins fonctionnels du système à développer ont été formalisés comme suit~:

\begin{itemize}
  \item \textbf{Interface Fournisseur (Espace FSI)}~: permettre au FSI d'ouvrir un ticket, de suivre son état en temps réel, et d'exécuter les tests techniques requis.
  \item \textbf{Interface Agent TT}~: permettre à l'agent de consulter ses tickets affectés, de renseigner les résultats d'intervention et de clôturer les tickets traités.
  \item \textbf{Dashboard de Performance Principale}~: restituer les KPIs sous forme de visualisations interactives, avec filtrage multicritère et alertes automatiques.
  \item \textbf{Module IA -- Pareto Intelligent}~: analyser automatiquement les données et générer des rapports de diagnostic contextualisés.
\end{itemize}

\subsection{Besoins non-fonctionnels}

\begin{table}[h!]
\centering
\caption{Synthèse des exigences non-fonctionnelles du système}
\begin{tabularx}{\textwidth}{L{3.5cm} X}
\toprule
\textbf{Exigence} & \textbf{Description} \\
\midrule
Performance    & Temps de chargement des pages inférieur à 3 secondes ; actualisation des KPIs quasi temps réel. \\
Disponibilité  & Disponibilité du système 24h/24, 7j/7 avec un taux de disponibilité cible $\geq$ 99,5~\%. \\
Sécurité       & Authentification sécurisée par rôle (FSI, Agent TT, Directeur Opérations) ; chiffrement des données en transit. \\
Scalabilité    & Architecture capable de supporter la montée en charge lors des pics d'incidents réseau. \\
Maintenabilité & Code source documenté, architecture modulaire facilitant l'ajout de nouvelles fonctionnalités. \\
\bottomrule
\end{tabularx}
\label{tab:nonfunc}
\end{table}

\subsection{Contraintes techniques et organisationnelles}

Le développement du système est soumis aux contraintes techniques suivantes~: utilisation du framework Next.js pour le développement frontend, intégration des données issues du fichier \texttt{all-data.xlsx} fourni par Hexabyte, compatibilité avec les navigateurs web modernes (Chrome, Firefox, Edge), et déploiement sur la plateforme Vercel.

% ═════════════════════════════════════════════
%  CHAPITRE 2
% ═════════════════════════════════════════════
\chapter{État de l'Art et Cadre Théorique}

Ce chapitre présente les fondements théoriques et technologiques sur lesquels repose le système développé. Nous abordons successivement les concepts de pilotage de la performance organisationnelle, la méthode Pareto et son application intelligente, les systèmes de gestion des interventions, et enfin les technologies retenues pour l'implémentation.

\section{Pilotage de la Performance Organisationnelle}

\subsection{Concepts fondamentaux des indicateurs clés de performance (KPI)}

Un indicateur clé de performance (KPI -- Key Performance Indicator) est une mesure quantifiable permettant d'évaluer le degré d'atteinte d'un objectif stratégique ou opérationnel au sein d'une organisation. Dans le domaine des services télécom, les KPIs jouent un rôle central dans le pilotage des opérations.

Pour être véritablement utile, un KPI doit satisfaire aux critères SMART~: Spécifique, Mesurable, Atteignable, Réaliste et Temporellement défini.

\subsection{Tableaux de bord de performance}

Un tableau de bord (\textit{dashboard}) de performance est un outil de visualisation synthétique permettant de présenter, en un coup d'œil, l'ensemble des indicateurs pertinents pour le pilotage d'une activité. Dans la littérature académique et managériale, le tableau de bord prospectif (\textit{Balanced Scorecard}) de Kaplan et Norton (1992) constitue l'une des références fondatrices en matière de pilotage stratégique par les KPIs~\cite{kaplan1996}.

\subsection{Analyse des écarts et alertes automatiques}

L'analyse des écarts consiste à comparer la valeur réelle d'un indicateur à sa valeur cible préalablement définie, afin de détecter les dérives de performance et d'alerter les responsables lorsqu'un seuil critique est franchi.

\section{Méthode Pareto et Analyse des Dysfonctionnements}

\subsection{Principe 80/20 de Pareto}

La loi de Pareto postule que, dans de nombreux phénomènes, 80~\% des effets sont produits par 20~\% des causes. Appliquée au domaine de la qualité industrielle et des services, cette loi constitue le fondement de l'analyse Pareto, outil incontournable de la démarche d'amélioration continue (Kaizen, Lean, Six Sigma)~\cite{imai1986}.

\subsection{Application dans les systèmes de gestion de réclamations}

Dans le contexte de Tunisie Telecom, l'analyse Pareto est appliquée aux codes de résolution des interventions techniques (\texttt{resolution\_code}). L'analyse des 22~851 tickets traités révèle que six codes de résolution concentrent 81,6~\% du volume total des réclamations.

\subsection{Intelligence artificielle pour une analyse Pareto intelligente}

L'apport de l'intelligence artificielle à l'analyse Pareto réside dans la capacité à automatiser la génération d'insights contextualisés à partir des données brutes. Dans le système développé, le module Pareto intelligent exploite l'API Groq avec le modèle Llama~3.3~70B pour analyser les patterns des données et générer automatiquement un rapport de diagnostic rédigé en langage naturel.

\section{Technologies et Architecture Retenues}

\subsection{Next.js -- Framework de développement frontend}

Next.js est un framework React de développement web fullstack, développé et maintenu par Vercel~\cite{nextjs2025}. Il offre des fonctionnalités avancées telles que le rendu côté serveur (SSR), la génération statique (SSG), et les API Routes permettant de développer des endpoints backend directement au sein de l'application.

\subsection{Groq API et modèle Llama~3.3}

Groq est une infrastructure d'inférence IA haute performance permettant d'exécuter des modèles de langage à grande échelle (LLM) avec des latences extrêmement faibles~\cite{groq2025}. Le modèle Llama~3.3~70B, développé par Meta AI, est un modèle de langage open-source de dernière génération offrant des capacités avancées de compréhension et de génération de texte en plusieurs langues~\cite{meta2024}.

\subsection{Recharts -- Bibliothèque de visualisation de données}

Recharts est une bibliothèque de visualisation de données React basée sur D3.js~\cite{recharts2025}. Elle offre une large gamme de types de graphiques avec des capacités d'animation et d'interactivité parfaitement adaptées aux besoins d'un dashboard de performance.

% ═════════════════════════════════════════════
%  CHAPITRE 3
% ═════════════════════════════════════════════
\chapter{Analyse et Conception du Système}

Ce chapitre est consacré à l'analyse des besoins et à la conception détaillée du système de pilotage de la performance. Nous présentons successivement la modélisation UML, l'architecture applicative globale, la conception des indicateurs de performance, le schéma de base de données et la conception des interfaces utilisateur.

\section{Modélisation des Besoins}

\subsection{Identification des acteurs}

L'analyse des besoins a permis d'identifier trois acteurs principaux interagissant avec le système~:

\begin{enumerate}
  \item \textbf{Le Fournisseur FSI (Hexabyte)}~: acteur externe qui ouvre les tickets de réclamation, suit leur évolution et exécute les tests techniques.
  \item \textbf{L'Agent Tunisie Telecom (Technicien)}~: acteur interne responsable de la réalisation des interventions terrain, qui renseigne les résultats et clôture les tickets traités.
  \item \textbf{Le Responsable de Performance (Directeur des Opérations)}~: acteur décisionnel qui exploite le dashboard pour piloter la qualité de service et identifier les dérives.
\end{enumerate}

\subsection{Architecture applicative globale}

Le système est articulé autour d'une architecture applicative en trois couches distinctes. La couche de présentation regroupe les trois interfaces utilisateur développées en Next.js et déployées sur Vercel. La couche métier est assurée par les API Routes Next.js qui implémentent la logique de calcul des KPIs. La couche données est constituée du fichier Excel \texttt{all-data.xlsx} fourni par Hexabyte, importé et traité côté serveur par la bibliothèque SheetJS.

\subsection{Description du diagramme de flux du système}

Le diagramme d'architecture illustre le flux complet de traitement d'une réclamation. Lorsqu'un client contacte le FSI pour signaler une anomalie, le FSI accède à l'Espace Fournisseur et recherche le numéro de ligne concerné. Après reparation, des tests techniques sont exécutés~: si les tests sont concluants, le ticket est clôturé ; en cas d'échec, l'intervention est relancée.

\section{Conception des Indicateurs de Performance}

\subsection{Cartographie des KPIs retenus}

Sur la base du diagnostic et des besoins exprimés par la direction des opérations, cinq indicateurs clés de performance ont été retenus et formalisés.

\begin{table}[h!]
\centering
\caption{Cartographie des KPIs retenus et leurs valeurs cibles}
\begin{tabularx}{\textwidth}{L{3.5cm} L{4cm} L{3cm} L{2.5cm}}
\toprule
\textbf{KPI} & \textbf{Définition} & \textbf{Formule} & \textbf{Valeur cible} \\
\midrule
Total Incidents    & Volume total de réclamations     & COUNT(tickets)                  & Suivi continu \\
Taux de Résolution & Part des tickets clôturés        & Clôturés / Total $\times$ 100   & $\geq$ 99,5~\% \\
Délai Moyen        & Durée moyenne de traitement      & AVG(TT\_resolution\_duration)   & $\leq$ 72~h (SLA) \\
Tickets Urgents    & Tickets avec délai $>$ 72~h      & COUNT(durée $>$ 72h)            & $<$ 5~\% du total \\
Délai Maximum      & Délai maximum observé            & MAX(TT\_resolution\_duration)   & Alerte si $>$ 240~h \\
\bottomrule
\end{tabularx}
\label{tab:kpis}
\end{table}

\subsection{Définition des seuils et alertes automatiques}

Chaque KPI est associé à un système d'alertes à trois niveaux~: vert (performance conforme à la cible), orange (dérive détectée -- attention requise) et rouge (seuil critique dépassé -- action corrective urgente). Sur les données réelles analysées, 23,1~\% des tickets déclenchent cette alerte critique.

\section{Conception de la Base de Données}

\subsection{Structure du fichier de données source}

La base de données opérationnelle du système est constituée du fichier \texttt{all-data.xlsx} fourni par le fournisseur Hexabyte. Ce fichier regroupe l'ensemble des tickets de réclamation traités sur l'année 2025, soit 22~851 enregistrements répartis en 17 champs.

\begin{table}[h!]
\centering
\caption{Dictionnaire des données de la base source (all-data.xlsx -- Hexabyte 2025)}
\begin{tabularx}{\textwidth}{L{4cm} L{2.5cm} X}
\toprule
\textbf{Champ} & \textbf{Type} & \textbf{Description} \\
\midrule
DLigne                  & Texte      & Identifiant unique de la ligne cliente \\
reference\_no Extranet  & Texte      & Référence du ticket sur le portail Extranet FSI \\
downtime\_start         & Date/Heure & Date et heure de début de l'incident signalé par le FSI \\
open\_time (HPSM)       & Date/Heure & Date et heure d'ouverture du ticket dans HPSM \\
Ref HPSM                & Texte      & Référence unique du ticket dans le système HPSM \\
problem\_status         & Texte      & Statut du ticket~: Clôturé (99,9~\%) ou Ouvert (0,1~\%) \\
close\_time             & Date/Heure & Date et heure de clôture du ticket \\
resolution\_code        & Texte      & Code de résolution de l'intervention TT \\
Region                  & Texte      & Région géographique de la ligne cliente (24 régions) \\
Offre                   & Texte      & Offre commerciale souscrite (ADSL Hybride, VDSL, etc.) \\
brief\_description      & Texte      & Description courte du motif de réclamation \\
product\_type           & Texte      & Famille d'intervention TT \\
TT\_resolution\_duration & Durée    & Durée totale de traitement du ticket (JJ HH:MM:SS) \\
\bottomrule
\end{tabularx}
\label{tab:dict}
\end{table}

\section{Conception des Interfaces Utilisateur}

\subsection{Vue globale du Dashboard Principale}

La vue globale constitue la page d'accueil du Dashboard Principale, accessible au Directeur des Opérations après authentification. Elle présente une synthèse macroscopique de la performance à travers trois KPI cards (Total Incidents~: 11~088, Taux de Résolution~: 99,9~\%, Délai Moyen~: 98,1 heures), complétées par deux graphiques~: la tendance mensuelle des interventions et le statut global des tickets sous forme de donut chart.

\subsection{Vue Réclamations -- Analyse approfondie}

La vue Réclamations propose une analyse détaillée des motifs de réclamation. Elle présente le Top~3 des problèmes par fréquence d'occurrence (Pas de synchronisation~: 46,5~\%, Modem synchro et pas d'accès internet~: 21,1~\%, Déconnexion courante~: 19,3~\%).

\subsection{Vue Interventions}

La vue Interventions présente l'analyse opérationnelle des 11~088 interventions du périmètre filtré. Un graphique en donut affiche la répartition par famille d'intervention~: Rien à signaler~: 32~\%, Fixation Coffret~: 26~\%, Remplacement Jarretière~: 21~\%, Relevé Isolement~: 14~\%, Remplacement câble~: 8~\%.

\subsection{Vue Performance -- Efficacité du Réseau \& Service}

La vue Performance est dédiée à l'analyse des délais de résolution selon différents axes. Elle présente trois métriques de temps (Temps moyen~: 4,1 jours, Temps maximum~: 181,0 jours avec alerte seuil critique, Temps minimum~: 0,0 jour), une carte thermique horizontale des délais moyens par région.

\subsection{Vue Pareto -- Analyse Pareto Intelligente}

La vue Pareto constitue le module analytique central du dashboard. Elle intègre un widget de génération de rapport Pareto intelligent, quatre KPI cards de synthèse, un graphique Pareto interactif des codes de résolution, et un tableau détaillé affichant pour chaque code de résolution~: volume, fréquence (\%), cumul Pareto (\%) et conformité SLA.

% ═════════════════════════════════════════════
%  CHAPITRE 4
% ═════════════════════════════════════════════
\chapter{Réalisation et Implémentation}

Ce chapitre présente en détail les développements réalisés pour chacune des composantes du système. Nous décrivons l'environnement technique, puis l'implémentation de chaque interface et module, en nous appuyant sur les données réelles analysées.

\section{Environnement de Développement}

\subsection{Stack technique utilisé}

Le système de pilotage de la performance organisationnelle a été développé en s'appuyant sur un ensemble de technologies modernes et éprouvées.

\begin{table}[h!]
\centering
\caption{Stack technique du système de pilotage de la performance}
\begin{tabularx}{\textwidth}{L{3.5cm} L{4cm} X}
\toprule
\textbf{Composant} & \textbf{Technologie} & \textbf{Rôle dans le système} \\
\midrule
Framework Web  & Next.js 14           & Développement frontend et API Routes backend \\
Langage        & TypeScript / JS      & Typage statique et développement applicatif \\
Visualisation  & Recharts             & Graphiques KPI, Pareto, tendances, camemberts \\
IA / LLM       & Groq API + Llama 3.3 & Génération des rapports Pareto intelligents \\
Données        & SheetJS (xlsx)       & Lecture et traitement du fichier all-data.xlsx \\
Style          & Tailwind CSS         & Design system et responsive layout \\
Déploiement    & Vercel               & Hébergement cloud et CI/CD automatisé \\
Versioning     & Git / GitHub         & Gestion du code source et collaboration \\
\bottomrule
\end{tabularx}
\label{tab:stack}
\end{table}

\subsection{Méthodologie de développement}

Le développement a été conduit selon une approche itérative et incrémentale, s'inspirant des principes de la méthodologie Agile. Chaque itération (sprint de deux semaines) a permis de livrer une fonctionnalité opérationnelle testable.

\section{Traitement des Données et Calcul des KPIs}

\subsection{Chargement et parsing du fichier source}

Le fichier \texttt{all-data.xlsx} contient 22~851 enregistrements de réclamations et interventions sur l'année 2025. Il est chargé côté serveur par la bibliothèque SheetJS et parsé en mémoire sous forme de tableau d'objets JavaScript. Chaque enregistrement est enrichi de champs calculés~: le délai de résolution en heures, le mois d'ouverture du ticket, le cluster d'offre (ADSL / VDSL / AUTRES-DATA) et le statut d'urgence.

\subsection{Calcul des indicateurs clés de performance}

Les cinq KPIs principaux sont calculés dynamiquement à chaque appel API. Sur le jeu de données complet (sans filtre), les valeurs observées sont~: Total Incidents~=~22~851, Taux de Résolution~=~99,9~\%, Délai Moyen~=~79,0 heures, Tickets Urgents~=~5~278 (23,1~\%), Délai Maximum~=~4~344 heures (181 jours).

\section{Implémentation du Dashboard Principale}

\subsection{Architecture des composants React}

Le Dashboard Principale est structuré en un ensemble de composants React modulaires et réutilisables. La navigation latérale (sidebar) donne accès aux sept vues principales~: Vue globale, Réclamations, Interventions, Performance, Offres \& Produits, Pareto, et les opérations détaillées.

\subsection{Vue Globale -- Tableau de bord exécutif}

L'histogramme de tendance mensuelle révèle une saisonnalité notable du volume d'incidents~: le pic observé en juillet 2025 (2~129 tickets) correspond à une période de forte chaleur impactant l'infrastructure réseau, tandis que les mois de février et mars affichent les volumes les plus faibles (1~708 tickets chacun).

\subsection{Vue Interventions -- Analyse opérationnelle}

Le graphique en donut révèle que la famille "Rien à signaler" représente 32~\% des interventions (4~631 tickets), ce qui signifie qu'un tiers des déplacements terrain ne débouche sur aucune action corrective côté Tunisie Telecom.

\subsection{Vue Performance -- Analyse des délais}

Trois axes d'analyse sont proposés~: les délais par région géographique (Mannouba~: 180,8~h, Sousse~: 160,0~h, Tataouine~: 145,8~h), les délais par cluster d'offre (ADSL~: 82,6~h, VDSL~: 67,2~h, Autres/Data~: 71,8~h), et les délais par famille d'intervention TT.

\subsection{Vue Régions -- Analyse territoriale}

La vue Régions met en évidence que la région de GAFSA concentre le volume le plus élevé (2~005 tickets), tandis que la région de Tataouine affiche le délai de traitement le plus long (10 jours de délai moyen).

\subsection{Vue Résolutions TT -- Analyse qualité}

La vue Résolutions TT présente l'analyse détaillée des 12 codes de résolution distincts. Le code principal est PC (Problème Câblage) avec 2~560 occurrences (23,1~\%), suivi d'Aucun avec 2~452 occurrences (22,1~\%).

\subsection{Vue Actions FSI -- Pilotage fournisseur}

La vue Actions FSI présente les 12 tickets actuellement ouverts (backlog à traiter), identifie l'action FSI prioritaire (Technique Local~: 5~042 cas), et liste les 11~088 actions qualifiées réparties par famille d'intervention.

\section{Implémentation du Module Pareto Intelligent}

\subsection{Architecture du module IA}

Le module Pareto intelligent est implémenté sous la forme d'une API Route Next.js qui effectue les opérations suivantes~: agrégation des données, calcul du tableau Pareto, construction d'un prompt structuré en français, appel à l'API Groq avec transmission du prompt, et restitution du rapport généré dans l'interface.

\subsection{Conception du prompt et calibrage du modèle}

La qualité du rapport généré repose en grande partie sur la conception du prompt soumis au modèle Llama~3.3~70B. Ce prompt est structuré en plusieurs sections~: un contexte opérationnel, les données Pareto agrégées, et une instruction de génération précisant le format attendu (rapport professionnel en français).

\subsection{Extrait du code source -- Module IA}

\begin{lstlisting}[style=jsStyle, caption={Appel API Groq pour le rapport Pareto intelligent}, label={lst:groq}]
// Construction du prompt et appel API Groq/Llama 3.3 70B
const prompt = `Analyse les donnees Pareto pour Tunisie Telecom :
  Periode: ${scope} | Total: ${total} reclamations
  Top causes: ${JSON.stringify(paretoData.slice(0,6))}`;

const response = await groq.chat.completions.create({
  model: 'llama-3.3-70b-versatile',
  messages: [{ role: 'user', content: prompt }]
});
\end{lstlisting}

\subsection{Résultats et restitution dans le dashboard}

Le rapport Pareto généré est restitué dans la vue Pareto du dashboard. Il présente une synthèse exécutive des performances de la période, une analyse détaillée des causes prioritaires (codes PC, Aucun et IPMSAN qui concentrent à eux seuls 63~\% des réclamations), des recommandations d'actions correctives, et une conclusion orientée vers l'amélioration continue.

% ═════════════════════════════════════════════
%  CHAPITRE 5
% ═════════════════════════════════════════════
\chapter{Tests, Validation et Résultats}

Ce chapitre présente la stratégie de tests mise en \oe uvre pour valider le système développé, les résultats obtenus pour chacun des scénarios de test, ainsi qu'une analyse critique des performances et des apports du système au regard des objectifs initiaux fixés.

\section{Stratégie de Tests}

\subsection{Types de tests réalisés}

La validation du système a été conduite à travers trois niveaux de tests complémentaires~:

\begin{itemize}
  \item \textbf{Tests unitaires}~: validation du comportement de chaque fonction de calcul de KPI isolément, en comparant les résultats obtenus avec les valeurs calculées manuellement à partir du fichier source.
  \item \textbf{Tests d'intégration}~: validation de la cohérence des données entre les différentes vues du dashboard. Un même périmètre de filtrage appliqué successivement sur différentes vues doit produire des totaux cohérents.
  \item \textbf{Tests fonctionnels}~: validation des scénarios utilisateur complets, de la sélection des filtres jusqu'à l'affichage des résultats et à la génération du rapport Pareto IA.
\end{itemize}

\subsection{Jeu de données de test}

Les tests ont été conduits sur le jeu de données réelles fourni par Hexabyte, constitué de 22~851 enregistrements couvrant l'intégralité de l'année 2025. La richesse et la variété du jeu de données (24 régions, 24 offres, 12 codes de résolution distincts, plus de 60 familles d'intervention, amplitude temporelle de 12 mois) constituent une garantie de robustesse pour le système.

\section{Résultats des Tests Fonctionnels}

\subsection{Validation du calcul des KPIs principaux}

\begin{table}[h!]
\centering
\caption{Validation des KPIs principaux}
\begin{tabularx}{\textwidth}{L{4cm} L{3cm} L{3cm} L{2cm} L{2cm}}
\toprule
\textbf{KPI} & \textbf{Réf. (Python)} & \textbf{Système} & \textbf{Écart} & \textbf{Statut} \\
\midrule
Total Incidents        & 22~851       & 22~851      & 0           & Conforme \\
Taux de Résolution     & 99,93~\%     & 99,9~\%     & $<$ 0,1~\%  & Conforme \\
Délai Moyen            & 79,0~h       & 98,1~h*     & Filtre 30j  & Conforme \\
Tickets Urgents        & 5~278 (23,1~\%) & 2~980 (27,0~\%)* & Filtre 30j & Conforme \\
Délai Maximum          & 4~344,4~h    & 4~344,4~h   & 0           & Conforme \\
\bottomrule
\multicolumn{5}{l}{\small * valeurs calculées sur le périmètre filtré 30 derniers jours}
\end{tabularx}
\label{tab:kpi-valid}
\end{table}

\subsection{Validation de l'analyse Pareto}

\begin{table}[h!]
\centering
\caption{Validation de l'analyse Pareto -- comparaison valeurs de référence / valeurs système}
\begin{tabularx}{\textwidth}{L{4cm} C{2cm} C{2.5cm} C{2cm} L{2.5cm}}
\toprule
\textbf{Code Résolution} & \textbf{Volume} & \textbf{Fréquence~\%} & \textbf{Cumul~\%} & \textbf{Statut} \\
\midrule
Aucun               & 5~665 & 24,8~\% & 24,8~\% & Conforme \\
PC                  & 5~119 & 22,4~\% & 47,2~\% & Conforme \\
IPMSAN              & 3~610 & 15,8~\% & 63,0~\% & Conforme \\
Distribution        & 1~913 &  8,4~\% & 71,4~\% & Conforme \\
Installation Interne& 1~213 &  5,3~\% & 76,7~\% & Conforme \\
Problème FSI        & 1~124 &  4,9~\% & 81,6~\% & Conforme \\
Câble de branchement& 1~112 &  4,9~\% & 86,5~\% & Conforme \\
Modem               &   730 &  3,2~\% & 89,7~\% & Conforme \\
\bottomrule
\end{tabularx}
\label{tab:pareto-valid}
\end{table}

\subsection{Validation du module Pareto Intelligent (IA)}

La validation du module Pareto intelligent a porté sur deux dimensions~: la pertinence des analyses générées et la cohérence des recommandations formulées. Une grille d'évaluation a été appliquée selon quatre critères~: exactitude factuelle, pertinence contextuelle, actionnabilité des recommandations, et qualité rédactionnelle.

Sur un ensemble de dix rapports générés en conditions variées, l'ensemble des critères d'évaluation ont été satisfaits.

\section{Présentation des Résultats Obtenus}

\subsection{Métriques de performance du système}

\begin{table}[h!]
\centering
\caption{Métriques de performance technique du système}
\begin{tabularx}{\textwidth}{L{5.5cm} L{4cm} L{4cm}}
\toprule
\textbf{Métrique} & \textbf{Valeur observée} & \textbf{Objectif cible} \\
\midrule
Temps de chargement Vue Globale   & $<$ 1,5 secondes  & $<$ 3 secondes \\
Temps de chargement Vue Pareto    & $<$ 2,0 secondes  & $<$ 3 secondes \\
Temps de génération rapport IA    & 3 à 8 secondes    & $<$ 15 secondes \\
Volume de données traité          & 22~851 enreg.     & $\geq$ 10~000 lignes \\
Disponibilité (Vercel)            & 99,99~\%          & $\geq$ 99,5~\% \\
Compatibilité navigateurs         & Chrome, Firefox, Edge, Safari & Navigateurs modernes \\
\bottomrule
\end{tabularx}
\label{tab:perf}
\end{table}

\subsection{Principaux enseignements tirés de l'analyse des données réelles}

L'analyse systématique des 22~851 tickets du jeu de données réel Hexabyte 2025 a permis de dégager plusieurs enseignements opérationnels majeurs~:

\begin{itemize}
  \item \textbf{Paradoxe performance/délais}~: le taux de résolution excellent (99,93~\%) coexiste avec un délai moyen élevé (79,0 heures) et une médiane bien plus faible (25,8 heures).
  \item \textbf{Concentration des causes}~: six codes de résolution concentrent 81,6~\% du volume des réclamations sur l'année.
  \item \textbf{Disparités régionales significatives}~: l'écart entre la région de Mannouba (délai moyen~: 180,8 heures) et les régions les plus performantes (délai $<$ 50 heures) révèle des inégalités importantes.
  \item \textbf{Dominance de l'ADSL}~: avec 66,4~\% du volume total (15~163 tickets), l'Offre de Gros ADSL Hybride représente le segment le plus exposé aux incidents.
  \item \textbf{Saisonnalité des incidents}~: le pic de juillet 2025 (2~129 tickets) et le creux de février 2025 (1~708 tickets) suggèrent une saisonnalité thermique des incidents réseau.
\end{itemize}

\subsection{Retour d'expérience et limites identifiées}

Malgré les résultats satisfaisants obtenus, plusieurs limitations du système actuel ont été identifiées~:

\begin{itemize}
  \item \textbf{Source de données statique}~: le système actuel s'appuie sur un fichier Excel sans connexion directe aux systèmes HPSM et Extranet, limitant l'aspect temps réel.
  \item \textbf{Périmètre fournisseur unique}~: les données disponibles ne couvrent que le fournisseur Hexabyte.
  \item \textbf{Absence de mécanisme d'authentification avancé}~: la version actuelle implémente une authentification basique par rôle.
  \item \textbf{Module IA dépendant d'un service tiers}~: la génération des rapports repose sur l'API Groq, service cloud externe.
\end{itemize}

% ═════════════════════════════════════════════
%  CONCLUSION GÉNÉRALE ET PERSPECTIVES
% ═════════════════════════════════════════════
\chapter*{Conclusion Générale et Perspectives}
\addcontentsline{toc}{chapter}{Conclusion Générale et Perspectives}

\section*{Synthèse des travaux réalisés}

Ce projet de fin d'année avait pour ambition de répondre à un besoin opérationnel concret et stratégique~: doter la direction des opérations de Tunisie Telecom d'un système de pilotage de la performance organisationnelle des interventions techniques. Pour atteindre cet objectif, nous avons conduit une démarche structurée en quatre phases complémentaires~: diagnostic, conception, implémentation et validation.

\section*{Objectifs atteints}

\begin{table}[h!]
\centering
\caption{Tableau récapitulatif des objectifs et livrables du projet}
\begin{tabularx}{\textwidth}{X L{4.5cm} C{2cm}}
\toprule
\textbf{Objectif initial} & \textbf{Livrable associé} & \textbf{Statut} \\
\midrule
Développer un système de pilotage de la performance & Dashboard Principale (7 vues) & Atteint \\
Intégrer une analyse Pareto intelligente (IA)       & Module Pareto + API Groq/Llama 3.3 & Atteint \\
Fournir un outil d'aide à la décision opérationnel  & Application déployée sur Vercel & Atteint \\
Concevoir une architecture multi-interfaces          & Interfaces FSI, Agent TT, Dashboard & Atteint \\
\bottomrule
\end{tabularx}
\label{tab:objectifs}
\end{table}

\section*{Limites et perspectives}

Conscients des limitations du système actuel, nous proposons plusieurs axes d'évolution pour les travaux futurs~:

\begin{itemize}
  \item Connexion directe aux systèmes HPSM et Extranet de Tunisie Telecom via des API REST dédiées, permettant un vrai pilotage en temps réel.
  \item Extension du périmètre à l'ensemble des fournisseurs FSI partenaires de Tunisie Telecom (au-delà d'Hexabyte).
  \item Développement d'un module de prévision basé sur des algorithmes de Machine Learning (régression, séries temporelles).
  \item Intégration d'une application mobile compagnon destinée aux agents techniciens.
  \item Déploiement d'un modèle LLM hébergé en local (\textit{on-premise}) au sein de l'infrastructure de Tunisie Telecom, afin d'éliminer la dépendance au service cloud Groq.
\end{itemize}

Ce projet constitue, à notre sens, une première pierre dans la construction d'un écosystème de pilotage intelligent des opérations télécom. Il démontre qu'il est possible, avec des technologies modernes et accessibles, de transformer des données opérationnelles brutes en instruments de décision puissants et actionnables.

% ═════════════════════════════════════════════
%  BIBLIOGRAPHIE
% ═════════════════════════════════════════════
\begin{thebibliography}{99}
\addcontentsline{toc}{chapter}{Bibliographie}

\bibitem{kaplan1996}
R.~S. Kaplan et D.~P. Norton, \textit{The Balanced Scorecard: Translating Strategy into Action}, Harvard Business School Press, Boston, 1996.

\bibitem{imai1986}
M.~Imai, \textit{Kaizen: The Key to Japan's Competitive Success}, McGraw-Hill, New York, 1986.

\bibitem{itil2019}
AXELOS, \textit{ITIL Foundation: ITIL 4 Edition}, The Stationery Office (TSO), Norwich, 2019.

\bibitem{iso2018}
ISO/IEC 20000-1:2018, \textit{Information Technology -- Service Management -- Part 1: Service Management System Requirements}, International Organization for Standardization, Geneva, 2018.

\bibitem{nextjs2025}
Vercel Inc., \textit{Next.js Documentation -- App Router}, Documentation officielle Next.js, \url{https://nextjs.org/docs}, consulté en 2025.

\bibitem{meta2024}
Meta AI Research, \textit{Llama 3: Open Foundation and Fine-Tuned Chat Models}, Meta AI Technical Report, 2024.

\bibitem{groq2025}
Groq Inc., \textit{Groq API Documentation -- Language Processing Units (LPU) Inference}, Documentation officielle Groq, \url{https://console.groq.com/docs}, consulté en 2025.

\bibitem{recharts2025}
Recharts Team, \textit{Recharts -- Redefined chart library built with React and D3}, Documentation officielle, \url{https://recharts.org}, consulté en 2025.

\bibitem{rockart1979}
J.~F. Rockart, ``Chief Executives Define Their Own Data Needs,'' \textit{Harvard Business Review}, vol.~57, no.~2, pp.~81--93, 1979.

\bibitem{deming1982}
W.~E. Deming, \textit{Out of the Crisis}, MIT Press, Cambridge, 1982.

\bibitem{tailwind2025}
Tailwind CSS Team, \textit{Tailwind CSS -- A Utility-First CSS Framework}, Documentation officielle, \url{https://tailwindcss.com/docs}, consulté en 2025.

\bibitem{sheetjs2025}
SheetJS Community, \textit{SheetJS CE -- Spreadsheet Data Toolkit}, Documentation officielle, \url{https://sheetjs.com/docs}, consulté en 2025.

\bibitem{tt2024}
Tunisie Telecom, \textit{Rapport Annuel 2024}, Direction de la Communication, Tunis, 2024.

\bibitem{ins2024}
Institut National de la Statistique (INS) Tunisie, \textit{Indicateurs du secteur des télécommunications en Tunisie}, INS, Tunis, 2024.

\end{thebibliography}

% ═════════════════════════════════════════════
%  ANNEXES
% ═════════════════════════════════════════════
\appendix
\chapter{Extraits du Code Source}

\section{Fonction de calcul des KPIs}

\begin{lstlisting}[style=jsStyle, caption={Fonction calculateKPIs -- Next.js API Route}, label={lst:kpi}]
// Calcul des KPIs principaux a partir du dataset filtre
function calculateKPIs(data) {
  const total = data.length;
  const closed = data.filter(
    t => t.problem_status === 'Cloture'
  ).length;
  const durations = data.map(
    t => parseDurationHours(t.TT_resolution_duration)
  );
  const urgent = data.filter(
    t => parseDurationHours(t.TT_resolution_duration) > 72
  ).length;
  return {
    totalIncidents:  total,
    resolutionRate:  (closed / total * 100).toFixed(1) + '%',
    avgDelay:        average(durations).toFixed(1) + 'h',
    urgentTickets:   urgent   // delai > 72h = SLA breach
  };
}
\end{lstlisting}

\section{Génération du Rapport Pareto IA}

\begin{lstlisting}[style=jsStyle, caption={Appel API Groq -- rapport Pareto intelligent}, label={lst:pareto}]
// Construction du prompt et appel API Groq/Llama 3.3 70B
const prompt = `Analyse les donnees Pareto pour Tunisie Telecom :
  Periode: ${scope} | Total: ${total} reclamations
  Top causes: ${JSON.stringify(paretoData.slice(0,6))}`;

const response = await groq.chat.completions.create({
  model: 'llama-3.3-70b-versatile',
  messages: [{ role: 'user', content: prompt }]
});
\end{lstlisting}

\chapter{Tableau Récapitulatif des Tests Fonctionnels}

\begin{table}[h!]
\centering
\caption{Tableau récapitulatif complet des tests fonctionnels -- Résultats de validation}
\begin{tabularx}{\textwidth}{L{3cm} L{4cm} L{4cm} C{1.5cm}}
\toprule
\textbf{Interface / Module} & \textbf{Scénario testé} & \textbf{Résultat attendu} & \textbf{Statut} \\
\midrule
Vue Globale        & Chargement sans filtre          & 22~851 tickets, 99,9~\%, 79h          & OK \\
Vue Globale        & Application filtre région GAFSA & 4~678 tickets, cohérence KPIs          & OK \\
Vue Interventions  & Répartition par famille         & Rien à signaler = 32~\%, Fixation = 26~\% & OK \\
Vue Performance    & Délai moyen par région          & Mannouba en tête (180,8~h)             & OK \\
Vue Pareto         & Graphique Pareto -- cumul 80~\% & Seuil 80~\% atteint à Problème FSI    & OK \\
Module IA          & Génération rapport Pareto mensuel & Rapport FR pro $<$ 10 secondes       & OK \\
Vue Régions        & Top régions par volume          & GAFSA = 2~005 tickets                 & OK \\
Vue Résolutions TT & Tableau Pareto codes résolution & PC = 2~560 (23,1~\%), cumul correct   & OK \\
Filtres globaux    & Changement de période           & Toutes les vues mises à jour           & OK \\
Responsive         & Affichage sur écran 1920×1080   & Layout correct, aucun débordement      & OK \\
\bottomrule
\end{tabularx}
\label{tab:tests}
\end{table}

\end{document}